% ===== §3 The Presuppositions of Traditional Eudaimonism =====
\section{The Presuppositions of Traditional Eudaimonism}
\label{sec:eudaimonia}

\noindent
The eudaimonist tradition is, in one sense, already more hospitable to the question of shared happiness than existentialism. Where existentialism begins with the individual thrown into a world of contingency and threat, eudaimonism begins with the question of the good life---of what it is for a human being to flourish---and this question has, from Aristotle onward, been understood to implicate the social and relational dimensions of human existence. Aristotle knew that human beings are political animals (\emph{z\={o}on politikon}), that friendship (\emph{philia}) is among the highest goods, that the virtues are exercised in a \emph{polis} and cannot be fully realised in isolation \citep{aristotle2000}. The eudaimonist tradition is not naive about relationality.

And yet it shares, at the level of its fundamental ontology, the same structural feature that limits the existentialist account: the subject of flourishing is the individual. \emph{Eudaimonia} is something that a person \emph{has}, or \emph{achieves}, or \emph{lives}. The others, the relations, the community, are conditions for flourishing or constituents of it, but they are conditions for and constituents of \emph{my} flourishing---the flourishing of a subject who is, in the last analysis, individual. This paper's task in the present section is to trace this structural feature through the major nodes of the eudaimonist tradition, showing not only where each account falls short of the shared flower but, more importantly, \emph{why} it falls short---what structural assumption would need to be revised for the account to accommodate what the shared flower reveals.

\subsection{Aristotle: \emph{Eudaimonia}, \emph{Phronesis}, and the Social Conditions of Individual Flourishing}
\label{subsec:aristotle}

\noindent
Aristotle's account of \emph{eudaimonia} in the \emph{Nicomachean Ethics} is the founding document of the eudaimonist tradition, and it remains, twenty-four centuries later, philosophically indispensable \citep{aristotle2000}. \emph{Eudaimonia}---often translated as happiness, but better rendered as flourishing or living-well-and-doing-well---is not a feeling or a subjective state but an activity: specifically, the activity of the soul in accordance with virtue (\emph{aret\={e}}) and, if there are several virtues, in accordance with the best and most complete. It is not something that happens to a person but something that a person does---or, more precisely, something that a person \emph{is} in the doing: the virtuous person does not merely act virtuously but acts from virtue, with the right motivation, at the right time, toward the right objects, in the right manner.

The social and relational dimensions of this account are real and important. The virtues---courage, justice, generosity, practical wisdom---are exercised in relation to others, and some of them, like justice and friendship-virtue, are constitutively relational: one cannot be just or a good friend in isolation. \emph{Phronesis}---practical wisdom, the capacity to deliberate well about what conduces to flourishing in particular circumstances---requires experience of the social world and is exercised in it. And \emph{philia}---friendship in Aristotle's rich sense, encompassing everything from utility-friendship to the highest friendship of those who share virtue and love each other for what they are---is not merely a condition for flourishing but one of its constituents: the happy person needs friends, and not merely as instruments or as mirrors of the self, but as genuine others whose good one pursues for their own sake.

All of this is true, and this paper will draw on it in \xref{sec:hermeneutics} when it re-reads these concepts under the GRB framework. But the structural point stands: \emph{eudaimonia} is \emph{my} activity, \emph{my} excellent functioning, \emph{my} flourishing. When Aristotle says that the happy person needs friends, he means that friendship is a constituent of \emph{my} happy life---that without friends, \emph{I} cannot fully flourish. The friend is not a co-subject of a shared flourishing; the friend is a condition for and partial constituent of \emph{my individual} flourishing. Even the highest \emph{philia}---in which two friends love each other for what they are, share a life, and pursue virtue together---is, in Aristotle's account, a relationship that constitutes and enriches two individual \emph{eudaimonia}s, not a third thing that is the flourishing of the relation itself.

\begin{claim}[The Aristotelian presupposition]
For Aristotle, the social and relational dimensions of flourishing are real and philosophically central. But they are understood as conditions for, and constituents of, the individual's \emph{eudaimonia}. The relational field is not itself the subject of flourishing; it is the medium in which individual flourishing is achieved and expressed.
\end{claim}

\subsection{Epicurus and the Stoics: Retreat from Relation}
\label{subsec:epicurusstoics}

\noindent
The Hellenistic schools represent, in different ways, a retreat from Aristotle's social eudaimonism---a move toward an account of flourishing that is more independent of the social world and, crucially, more independent of what happens in it. This retreat is philosophically motivated: if flourishing depends on relations, communities, and external goods, then it is vulnerable to fortune, and the Hellenistic thinkers were deeply concerned with providing an account of happiness that was, to the greatest possible degree, within the agent's own power.

Epicurus locates \emph{eudaimonia} in pleasure (\emph{h\={e}don\={e}}), specifically in the absence of pain (\emph{aponia}) and mental disturbance (\emph{ataraxia}) \citep{epicurus1994}. But the pleasure that matters is not sensory excitement---which typically brings more pain in its wake than it is worth---but the calm pleasure of a life free from fear, desire, and disturbance. Epicurus does value friendship highly: the friends of the Garden are central to his account of the pleasant life, and his famous remark that it is more pleasant to give than to receive a benefit points toward something like a relational account of pleasure. But the fundamental structure remains individualist: \emph{ataraxia} is a state of \emph{my} soul, achieved by \emph{my} philosophical practice, in which \emph{I} am no longer disturbed by false beliefs about death, the gods, and the nature of pleasure. The friends of the Garden are conditions for this state, not co-subjects of a shared state.

The Stoics push further in the direction of self-sufficiency \citep{marcus2002}. For the Stoic sage, \emph{eudaimonia} consists in virtue alone, and virtue is entirely within the agent's power regardless of external circumstances. The preferred indifferents (\emph{proh\={e}gmena adiaphora})---health, wealth, friendship, reputation---are to be pursued when available but do not constitute happiness; only virtue does. This means that the Stoic sage can be happy on the rack, and that the loss of friends, community, or all external goods leaves happiness intact. The relational world is, for Stoicism, philosophically peripheral to flourishing: cosmopolitan concern for all rational beings (\emph{oikei\={o}sis}) extends one's moral concern to the whole of humanity, but this extension is an expression of the sage's virtue, not a dependence of the sage's happiness on any particular relation.

\emph{Apatheia}---the Stoic ideal of freedom from the passions---is, from the GRB perspective, the clearest formulation of what relational happiness must oppose. For the emotions that arise in the shared encounter with the flower---the joy, the resonance, the sense of being touched together by the contingent---are precisely the kind of affective engagement that \emph{apatheia} aims to transcend. The Stoic sage has achieved a happiness that is invulnerable to the flower's beauty and indifferent to whether it is shared---a happiness that is, by that very invulnerability, no longer capable of the relational co-evolution that the GRB framework will identify as flourishing's highest form.

\subsection{Csikszentmihalyi: Flow and the Individual in Optimal Experience}
\label{subsec:flow}

\noindent
The modern psychological tradition has given the eudaimonist inheritance new empirical content, and no concept has been more influential in this tradition than Csikszentmihalyi's account of flow \citep{csikszentmihalyi1990}. Flow is the state of optimal experience that occurs when a person is fully absorbed in a challenging activity, when skill and challenge are matched, when self-consciousness recedes and time is distorted, when the activity is intrinsically rewarding and the person feels a sense of effortless control. Athletes, musicians, surgeons, chess players, and rock climbers have all reported this state, and Csikszentmihalyi's research has documented its structure with considerable empirical precision.

The concept is philosophically significant: flow is not mere pleasure but engagement, not the satisfaction of desire but the full exercise of capacity, and in this sense it is closer to Aristotelian \emph{energeia} than to Epicurean \emph{h\={e}don\={e}}. The person in flow is not passive but active, not receiving but doing, and the doing is its own reward. This captures something important about the structure of happiness that the hedonic tradition misses.

And yet the structural presupposition is evident: flow is an individual's experience, arising from the individual's engagement with a task. The challenge-skill balance is the individual's challenge-skill balance. The absorption is the individual's absorption. Even when flow occurs in a social activity---team sports, musical ensemble, collaborative work---the flow is understood as something that each individual player or musician or worker either achieves or does not achieve. The team's victory may create the conditions for the individual player's flow; it does not itself constitute a relational flow that is irreducible to the sum of individual flows.

This is not a trivial limitation. Consider the scene with which this paper began: two people sitting together, doing nothing, in the kind of still and wordless co-presence that is one of the deepest forms of shared happiness. There is no task. There is no challenge. There is no skill being exercised. On Csikszentmihalyi's account, there is no flow---and therefore, on the most influential modern account of optimal experience, no happiness of the kind that concerns us here. The gap between what the theory predicts and what anyone who has experienced this kind of shared stillness knows to be true is itself a philosophical datum. It points to the need for a concept of relational flow that is not derivative of individual flow but genuinely different in kind---a concept that \xref{subsec:flow-grb} will develop.

\subsection{Maslow: Self-Actualisation and the Hierarchy of Individual Needs}
\label{subsec:maslow}

\noindent
Abraham Maslow's hierarchy of needs is perhaps the most widely known framework in the psychology of human motivation and flourishing \citep{maslow1954}. The hierarchy moves from the physiological needs at the base---food, water, shelter, sleep---through safety, belonging and love, and esteem, to self-actualisation at the apex: the need to realise one's full potential, to become what one is capable of becoming. Maslow's self-actualising persons are characterised by a rich set of qualities---peak experiences, democratic character, deep interpersonal relations, autonomy, continued freshness of appreciation---that give the concept philosophical substance beyond mere self-help.

The relational dimensions of the hierarchy are acknowledged: belonging and love constitute the third level, and Maslow's self-actualisers are typically described as having profound relationships with a few significant others. But the structure of the hierarchy is unmistakably individualist: the higher needs are the individual's needs, self-actualisation is the individual's actualisation, and the relational goods at level three are understood as needs whose satisfaction enables the individual to ascend toward the apex. The other is the satisfier of my belonging-need, the partner in my love, the mirror of my esteem---but the subject whose needs are being met, whose potential is being realised, whose peak experiences are being had, is always and only the individual.

Under GRB, this hierarchy requires not merely modification but inversion at its apex. If the subject of flourishing is the relational field, then what Maslow places at the top---individual self-actualisation, the realisation of the individual's full potential---is not the highest form of flourishing but a penultimate one. The highest form is the co-evolution of a relational system that generates something irreducible to either participant: not my self-actualisation, not your self-actualisation, but the relational actualisation of a shared field that neither of us could have become alone.

\subsection{Self-Determination Theory: Relation as Need, Not Ground}
\label{subsec:sdt}

\noindent
Deci and Ryan's self-determination theory (SDT) represents one of the most sophisticated and empirically grounded accounts of human motivation and flourishing in contemporary psychology \citep{deci1985}. SDT identifies three basic psychological needs whose satisfaction is necessary for well-being: autonomy (the experience of volition and self-endorsement of one's actions), competence (the experience of effectiveness and mastery), and relatedness (the experience of meaningful connection with others). All three needs are universal, and the failure to satisfy any one of them undermines well-being regardless of the satisfaction of the others.

The inclusion of relatedness as a basic need is philosophically significant, and SDT's empirical research on the conditions for relatedness---its findings on the role of autonomy-supportive environments, on the damage done by controlling relationships, on the importance of genuine rather than contingent regard---has genuine philosophical implications. SDT knows that relationships matter for flourishing and has studied how and why they matter with empirical rigour.

But the structural presupposition is once again evident: relatedness is a \emph{need} of the individual, a need whose satisfaction contributes to the \emph{individual's} well-being. The relational other is the satisfier of my relatedness need---or, in the more nuanced SDT account, the co-participant in an interaction that either supports or thwarts my autonomy, competence, and relatedness. The relational field is not itself a subject of flourishing; it is the medium in which individual needs are satisfied or frustrated.

\begin{claim}[The SDT presupposition]
Self-determination theory's inclusion of relatedness as a basic need marks a genuine advance over purely individualist accounts of well-being. But relatedness remains, in SDT, a need of the individual rather than an ontological ground of the subject. The theory asks how the individual's need for relatedness is satisfied; it does not ask whether the individual is, at a deeper level, constituted by and in the relational field.
\end{claim}

\subsection{Frankl: Meaning and the Individual Who Finds It}
\label{subsec:frankl}

\noindent
Viktor Frankl's account of meaning as the fundamental human motivation---forged in the extremity of the concentration camp and developed in his logotherapy---is philosophically and humanly one of the most powerful in the modern tradition \citep{frankl1959}. Frankl's insight that human beings can endure almost any \emph{how} if they have a \emph{why}---that meaning, not pleasure or power, is the primary human drive---has the kind of existential weight that only lived testimony can give to a philosophical claim.

For Frankl, meaning can be found in three ways: through what one gives to the world (creative values), through what one receives from the world (experiential values, including love and beauty), and through the attitude one takes toward unavoidable suffering (attitudinal values). The second category---experiential values, including love---points most directly toward relational happiness: the encounter with another person in love, for Frankl, is one of the highest forms of meaning-discovery.

But Frankl's account remains, structurally, an account of how the individual finds or creates meaning. Even in love, what happens is that \emph{I} encounter \emph{you} and, in that encounter, \emph{I} discover the full depth of your personhood---your uniqueness, your irreplaceability, the specific way in which you are you rather than anyone else. This discovery is profoundly significant, and Frankl describes it with genuine philosophical care. But the subject who makes the discovery, who finds meaning, who is constituted by the encounter, is still the individual: the encounter enriches me, transforms me, gives me meaning; it does not produce a third ontological unit---a relational being---whose flourishing is distinct from and irreducible to mine.

\subsection{Seligman's PERMA and the Aggregation of Individual Well-Being}
\label{subsec:perma}

\noindent
Martin Seligman's PERMA model---Positive emotion, Engagement, Relationships, Meaning, and Accomplishment---is the most influential framework in contemporary positive psychology \citep{seligman2011}. It is explicitly pluralist: well-being is not a single thing but a multidimensional construct, and no single component is sufficient or necessary for flourishing. The inclusion of Relationships as a distinct component marks a significant development from earlier hedonic models.

But the structure of aggregation that the PERMA model employs---well-being as the sum or combination of five components, each of which contributes to the individual's overall flourishing---reproduces, at the level of theory, the individualist presupposition. Relationships are good for \emph{my} well-being; positive emotions are \emph{my} positive emotions; engagement is \emph{my} engagement; meaning is the meaning \emph{I} find; accomplishments are \emph{my} accomplishments. The relational component (R) is not fundamentally different in kind from the others: it is one input among five into the individual's overall well-being score.

Under GRB, this aggregation misses something essential. The happiness of the shared flower is not my positive emotion (P) plus a contribution to my relationships (R): it is a qualitatively different phenomenon that cannot be decomposed into individual components because it does not arise from the individual but from the relational field. The PERMA model, for all its sophistication, is a model of the individual's well-being in a social context, not a model of the well-being of a relational system.

\subsection{The Shared Presupposition and What It Costs}
\label{subsec:eudaimdiagnosis}

\noindent
Across the eudaimonist tradition---from Aristotle's virtuous activity through the Stoics' self-sufficient sage, from Csikszentmihalyi's flow through Maslow's self-actualisation, from Frankl's meaning through Seligman's PERMA---a single structural presupposition runs: \emb{the subject of flourishing is the individual}. The others, the relations, the community, the shared encounters with contingent beauty, are conditions for, constituents of, or contributors to individual flourishing. They are not themselves the subjects of flourishing.

This presupposition is not arbitrary. It reflects a deep and in many respects accurate picture of human experience: I do experience my own happiness; I am the one who feels the joy, notices the flower, is moved by the encounter. There is something that it is like to \emph{be me} encountering the flower, and this something is genuinely mine in a way that cannot be entirely dissolved into the relational field. The paper does not deny this. What it denies is that this is the whole story---that the individual's experience of happiness exhausts the phenomenon, that the shared encounter with the flower can be fully accounted for by decomposing it into individual experiences and summing.

What the shared flower reveals is that there is a level of the phenomenon that the individualist presupposition systematically misses: the level of the relational field itself, in which something happens that is not in me and not in you but between us, and in which the happiness that arises is an emergent property of that field rather than a property of either of us. To account for this level, a different framework is needed---one that begins not with the individual who then enters into relations but with the relational field that generates individuals and in which they continue to be constituted. That framework is what the next section introduces.
